% ---------------------------------------------------------------------------
% Chapter 12: The Census Takers: Particle Flow
% Generated from the outline; edit freely, this file is now the source.
% ---------------------------------------------------------------------------
\chapter{The Census Takers: Particle Flow}
\chaptermark{The Census Takers}
\label{ch:census_takers}

\begin{journeybox}
The census takers walk through the whole detector with one rule: count everyone
once and only once. Tracks are linked to clusters, charged hadrons are separated
from photons and neutral hadrons, and the border guards' low readings are
corrected particle by particle.
\end{journeybox}

\physicsline{PF blocks, linking, charged hadrons, photons, neutral hadrons, PF hadron calibration, jet energy composition}

\section{Counting everyone once and only once: the Particle Flow philosophy}
\label{sec:census_takers:counting_everyone_once_and_only_once_the}
\todo{Write this section.}

\section{Linking tracks and clusters; the PF hadron calibration $a(E)$, $b(E)$}
\label{sec:census_takers:linking_tracks_and_clusters_the_pf_hadro}
\todo{Write this section.}

\section{The jet energy composition: charged hadrons, photons, neutral hadrons, leptons}
\label{sec:census_takers:the_jet_energy_composition_charged_hadro}
\todo{Write this section.}

\section{Why Particle Flow restores the response towards one}
\label{sec:census_takers:why_particle_flow_restores_the_response_}
\todo{Write this section.}

\section{Our jet at this stage: the list of PF candidates of our jet}
\label{sec:census_takers:our_jet_at_this_stage_the_list_of_pf_can}
\todo{Numbers, plots and event display of the $b$-jet at this stage.}

\begin{ledger}{The Census Takers}
  \ledgerrow{before this stage}{}{}{}{--}
  \ledgerrow{after this stage}{}{}{}{}
\end{ledger}

\section{Further reading}
\label{sec:census_takers:further_reading}
Official and theory references for this stage: \cite{CMS:2017yfk}.

\begin{pitfalls}
\todo{Pitfalls and FAQs for newcomers at this stage.}
\end{pitfalls}

\begin{exercises}
\item \todo{Exercise 1.}
\item \todo{Exercise 2.}
\end{exercises}
